\section{TYPE-004: Unvalidated JSON Parsed Results}
\label{sec:type-004}
\subsection*{Metadata}
\begin{tabular}{ll}
\textbf{Issue ID} & TYPE-004 \\
\textbf{Severity} & MEDIUM \\
\textbf{Category} & Type Safety \\
\textbf{Branch} & \branchref{fix/TYPE-004-validate-json-parse-in-task-tracker}
\textbf{Changes} & \branchdiff{fix/TYPE-004-validate-json-parse-in-task-tracker} \\
\end{tabular}

\subsection*{Executive Summary}
The task tracker at \srcfile{src/tools/task-tracker.ts:39-40} uses \code{JSON.parse(raw) as TasksStore} to load persisted task data. The \code{as TasksStore} cast tells TypeScript to trust that the raw JSON matches the \code{TasksStore} interface, but there is no runtime validation. If the file is corrupted, contains unexpected data, or was written by a different code version, the cast hides type mismatches and causes runtime errors downstream.

\subsection*{Root Cause Analysis}
The parsed result is cast without validation:

\begin{lstlisting}[language=TypeScript]
// BEFORE — unchecked cast
const raw = await readFile(tasksPath, "utf-8");
const tasksData = JSON.parse(raw) as TasksStore;  // No validation
\end{lstlisting}

\subsection*{Fix Implementation}
A type guard validates the parsed structure:

\begin{lstlisting}[language=TypeScript]
// AFTER — runtime type guard
function isTasksStore(data: unknown): data is TasksStore {
    if (typeof data !== "object" || data === null) return false;
    const d = data as Record<string, unknown>;
    return Array.isArray(d.tasks) && d.tasks.every(
        (t: unknown) =>
            typeof t === "object" && t !== null &&
            typeof (t as Record<string, unknown>).id === "string" &&
            typeof (t as Record<string, unknown>).objective === "string" &&
            Array.isArray((t as Record<string, unknown>).steps),
    );
}

async function loadTasks(gitagentDir: string): Promise<TasksStore> {
    const tasksFile = join(gitagentDir, "learning", "tasks.json");
    try {
        const raw = await readFile(tasksFile, "utf-8");
        const parsed = JSON.parse(raw);
        if (!isTasksStore(parsed)) {
            console.warn("Invalid tasks.json format, starting fresh");
            return { tasks: [] };
        }
        return parsed;
    } catch {
        return { tasks: [] };
    }
}
\end{lstlisting}

\subsection*{Affected Files}
\begin{itemize}
    \item \srcfile{src/tools/task-tracker.ts} — loadTasks() function
    \item \srcfile{src/\_\_tests\_\_/type-004.test.ts}
\end{itemize}

\subsection*{Verification}
Invalid JSON now returns an empty store instead of silently producing incorrect types.
